\chapter{La machine au complet}
% version complète: chapters/10_synthesis.tex

\pencilfig{10}{\pencilart{10}}{La machine au complet: le bus de données, les freins qu'il contient et, au-dessus, quatre canaux radio de contrôle.}

Nous avons pris les types de neurones un par un en demandant ce que chacun ajoute à la machine. Il est temps d'assembler les pièces.

On obtient une machine à trois couches. En bas, l'immense réseau téléphonique des neurones \nt{GLUT}: les données y circulent. À côté, les neurones inhibiteurs \nt{GABA}: ils gèrent le flux, décident de ce qui passe, quand et à quel volume. Et au-dessus de tout cela, quelques stations de radio, chacune avec son bouton. \nt{DOPA} dit quelles modifications des connexions garder. \nt{SERO} tient le niveau général et, selon l'hypothèse, l'horizon de patience. \nt{NORA} choisit entre chercher et décider. \nt{ACET}, entre écrire et lire. Chaque station est en réalité un petit faisceau de canaux plutôt qu'un fil unique, mais cela ne fait toujours qu'une poignée de canaux pour quatre-vingt-six milliards de neurones.

\section*{Une histoire}

Suivons un événement à travers toute la machine. L'animal sait depuis longtemps qu'en appuyant sur le levier A, il obtient du jus. Un jour, le monde change: A ne marche plus, et c'est le levier B, dont l'animal ne se servait presque jamais, qui donne le jus.

Le jus ne vient pas après A: la dopamine répond par un creux, et les connexions qui ont choisi A s'affaiblissent. Les creux se répètent, les erreurs deviennent plus grandes que d'habitude: la noradrénaline, selon l'hypothèse, signale que le monde a changé. Le bouton de contraste baisse, le choix devient plus souple, et le bruit pousse plus souvent à essayer B. L'acétylcholine fait passer le réseau en mode écriture. L'essai de B apporte un jus inattendu: la dopamine s'embrase, les connexions qui ont choisi B se renforcent. Au bout de quelques essais, l'attente rattrape le monde nouveau, les surprises cessent, les boutons reviennent en place.

Remarquez: aucun bouton ne sait ce que font les autres. Chacun répond à ses propres indices, et l'ordre des opérations se met en place tout seul, comme dans un circuit asynchrone où chaque bloc se déclenche quand ses entrées sont prêtes. Que chaque bouton se comporte ainsi pris isolément est, pour l'essentiel, mesuré. Qu'ils travaillent ensemble de façon coordonnée est encore une hypothèse.

\section*{En quoi elle diffère de votre ordinateur portable}

L'ordinateur portable a une horloge commune, le cerveau non. Les éléments de l'ordinateur sont fiables; une synapse perd plus de la moitié des signaux. La mémoire de l'ordinateur est séparée du processeur; dans le cerveau, elle loge dans les connexions mêmes. L'ordinateur apprend (s'il apprend) par rétropropagation de l'erreur; le cerveau, par des règles locales et un unique \og mieux que prévu\fg{} diffusé à tous. Et surtout: l'ordinateur dépense des dizaines de watts pour un seul processeur, le cerveau vingt watts pour tout.

\section*{Ce que nous ignorons}

Nous ne savons pas dans quelle mesure le cortex utilise l'instant précis des clics. Nous ne savons pas comment il règle des milliards de connexions dans des circuits à nombreux étages, alors qu'il n'a qu'un seul professeur pour tous. Nous ne comprenons pas entièrement ce que transmettent les stations de radio. Nous ne savons pas comment des parties éloignées du cerveau coordonnent leur travail sans horloge commune. Et personne ne sait encore construire une machine qui ferait la même chose avec vingt watts.

La suite du livre sort de ce noyau: vers les entrées et les sorties de la machine, son débogage, le sommeil et les machines faites de neurones vivants sur une puce.

\fullversion{10}
